\section{Discussion}
\label{sec:discussion}

We have proposed a spline-augmented Fay--Herriot estimator that relaxes the
linearity of the regression mean for a single continuous covariate while
retaining the simplicity of an area-level model. Because the spline is
represented as a second random effect, the model can be fitted with standard
REML machinery, and the smoothing parameter is selected automatically as part
of the variance component estimation. The parametric bootstrap provides MSE
estimates that account for the estimation of both variance components, and our
simulations indicate that the resulting intervals have close to nominal
coverage even with as few as 50 districts.

In the NPHNS application, the SAFH estimator changed the substantive
conclusions in a policy-relevant way: four additional districts were classified
as exceeding the national stunting target, all of them in the eastern
highlands, where the standard model underestimates prevalence. This illustrates
a general point about small area estimation that is easy to overlook. Model
misspecification does not merely add noise to the district estimates; it
produces bias that is correlated with district characteristics, and because
shrinkage is strongest where samples are smallest, the bias is concentrated in
the districts that are usually of greatest policy interest.

Several limitations should be noted. First, the SAFH model allows a nonlinear
effect for only one covariate. Additive models with several smooth terms are a
natural extension, but with a few hundred districts the data contain limited
information about several smooth functions simultaneously, and we expect the
REML estimates of the corresponding variance components to be unstable. Second,
like all Fay--Herriot type models, SAFH treats the smoothed sampling variances
$\psi_i$ as known. Errors in the generalized variance function propagate to the
weights $\hat w_i$, and a fully Bayesian treatment that models the sampling
variances jointly with the prevalences would be a useful extension
\citep{you2006small}. Third, the remoteness index is itself estimated from a
travel-time surface, and measurement error in $r_i$ could attenuate the
estimated function $\hat f$ \citep{ybarra2008small}.

Several limitations should be noted. First, the SAFH model allows a nonlinear
effect for only one covariate. Additive models with several smooth terms are a
natural extension, but with a few hundred districts the data contain limited
information about several smooth functions simultaneously, and we expect the
REML estimates of the corresponding variance components to be unstable. Second,
like all Fay--Herriot type models, SAFH treats the smoothed sampling variances
$\psi_i$ as known. Errors in the generalized variance function propagate to the
weights $\hat w_i$, and a fully Bayesian treatment that models the sampling
variances jointly with the prevalences would be a useful extension
\citep{you2006small}. Third, the remoteness index is itself estimated from a
travel-time surface, and measurement error in $r_i$ could attenuate the
estimated function $\hat f$ \citep{ybarra2008small}.

Finally, the estimated function $\hat f$ should not be given a casual
interpretation. Remoteness is correlated with many unobserved determinants of
child growth, including dietary diversity, disease burden and access to
antenatal care, and the SAFH model is designed for prediction rather than for
identifying the effect of travel time on stunting. Policies that reduce travel
times, such as road construction, may therefore have effects on nutrition that
differ substantially from what the slope of $\hat f$ would suggest.

The approach extends directly to other indicators routinely published from
household surveys, such as wasting, underweight or vaccination coverage, and to
other auxiliary variables with a plausibly nonlinear relationship to the
outcome, such as night-time light intensity or population density. An
implementation in R is provided in the supplementary material.
